%=== FRONT PART ===
%=== ABSTRCT ===

%\begin{center}
\chapter*{Abstract}
\rhead{Abstract}
%\end{center}
\addcontentsline{toc}{chapter}{Abstract}
%===============================
% ABSTRACT GOES HERE - 250 words limit
%===============================
As companies move toward containerisation technologies to improve scalability and portability in their applications, Kubernetes has emerged as the industry standard for orchestrating containerised workloads. Despite its advantages, deploying containerised workloads to the cloud is highly complex, and resource intensive. Developers and engineer also need to master the various cloud technologies, such as security and networking requirements, often diverting focus from core application development. 
This project introduces cloudkube, a lightweight, Python-based command-line tool designed to automate the deployment of Kubernetes clusters in the cloud. Inspired by minikube and leveraging Infrastructure as Code (IaC) principles with Terraform, cloudkube abstracts the complexities of cloud provisioning and Kubernetes setup. It enables developers to focus on core application development, rather than cloud or deployment configurations.. The tool supports customizable configurations, automated local environment setup, and idempotent resource provisioning while ensuring low-cost prototyping by leveraging AWS Free Tier resources.
Through cloudkube, deployment times are significantly reduced, and infrastructure provisioning is made accessible to users without deep expertise in cloud technologies. This paper details the tool’s architecture, implementation, and workflow, showcasing its potential to enhance productivity, reduce operational overhead, and accelerate the deployment of containerised applications to the cloud. 
%===============================
\\
\par
\textbf{Keywords:} Containerization, Cloud Automation, Infrastructure as Code (IaC), Terraform, Cloud Deployment, Kubernetes Orchestration, Developer Productivity, Lightweight CLI Tools, Cloud-Native Applications

%=== END OF ABSTRACT ===
